%!TEX root = ../template.tex %
% !TeX spellcheck = en_US
%%%%%%%%%%%%%%%%%%%%%%%%%%%%%%%%%%%%%%%%%%%%%%%%%%%%%%%%%%%%%%%%%%%%
%% appendix2.tex
%% NOVA thesis document file
%%
%% Chapter with example of appendix with a short dummy text
%%%%%%%%%%%%%%%%%%%%%%%%%%%%%%%%%%%%%%%%%%%%%%%%%%%%%%%%%%%%%%%%%%%%

\typeout{NT FILE appendix2.tex}%

\chapter{Support information to Chapter~\ref{ch:Optical_hBN_superlattice}}
\label{app:hBN_superlattice}


\section{Recovering the original Hamiltonian}
\label{sec:recover_original_hamiltonian}

To conciliate the renormalized low-energy Hamiltonian~\eqref{eq:new_Dirac_Hamiltonian} with the original one~\eqref{eq:low_energy_hamiltonian_hBN}, we take in the former the two limits $V_0 \to 0$ and $L \to \infty$ separately.
In both cases, we should recover the latter.
We take first the limit $V_0 \to 0$ to obtain
%
\begin{equation}
    \bar{H}_{m} = \frac{m}{2}G_{0}\hbar v_{F}+\hbar v_{F}\sigma_{x}q_{x}
\end{equation}
%
\noindent for $m \neq 0$, and
%
\begin{equation}
\bar{H} =\hbar v_{F}[\sigma_{x}q_{x}+\sigma_{y}q_{y}]-\sigma_{z}M
\end{equation}
%
\noindent for $m = 0$. 
The latter case is simply the original low-energy Hamiltonian of hBN in the absence of potential as expected.
There is only a difference in the sign of the mass term, but that is not critical.
The result for $m \neq 0$ is not null, but we can interpret the two terms on the \gls{rhs} as follows.
The first term  is purely a global shift in energies resultant from shifting the momentum along the $x$ direction by $m G_0/2$, a consequence of using basis states with the spatial dependent part of the form $\exp{\ii (\bq \pm m G_0 \hat{\mathbf{x}}/2)}$ used when deriving the renormalized low-energy Hamiltonian.
The second term comes from the measurement of momentum from the edges of the mini Brillouin zone along the $x$ direction, which no longer exists since $V_0 = 0$.
Therefore, it is not relevant.

We take now the limit $L \to \infty$, that is, the limit of a constant potential.
For each value of $m$ we have a Hamiltonian, but if we gather them all by summing over $m$ we get

\begin{align}
H_{\text{total}} &= 
\sum_{m} \overline{H}_{m} = 
\nonumber\\
&= \sum_{m} \frac{m}{2} G_{0} \hbar v_{F} +
\hbar v_{F} [\sigma_{x}q_{x}+\sigma_{y} q_{y} \sum_{m}J_{m}(\beta)] -
\sigma_{z} M \sum_{m} J_{m}(\beta) = 
\nonumber\\
& \overset{L\to \infty}{=} \hbar v_{F}
[\sigma_{x}q_{x}+\sigma_{y}q_{y}] - \sigma_{z}M,
\end{align}
%
\noindent where to arrive at the final result we used the fact that $G_0 \to 0$ when $L \to \infty$ and the rule $\sum_{m} J_{m}(\beta) = 1$, valid for the Bessel functions of the first kind.
In this way, we are able to recover the original low-energy Hamiltonian of hBN (with the small caveat of the sign in the mass term once more).

\section{Exchange energy correction to the band gap}
\label{sec:exchange_energy}

We recall the Bethe--Salpeter equation:

\begin{equation}
    \begin{split}
        E_\nu \frac{1}{\sqrt{A}} \phi_\nu(\bk)  =&   \phi_{\nu}(\bk) (E_{\bk,c}-E_{\bk,v}) - \frac{1}{A} \sum_{\bp} \phi_{\nu} (\bp) V_{\bk-\bp} F_{cvvc}(\bk,\bp,\bk-\bp) +  \\
        & + \frac{1}{A} \phi_{\nu}(\bk) \sum_{\bp} V_{\bp} \left[F_{vvvv}(\bp + \bk ,\bk,\bp) - F_{vcvc}(\bp + \bk,\bk,\bp) \right].
    \end{split}
\end{equation}

\noindent We wish to compute the contribution to the band gap due to the exchange energy coming from the third term in the \gls{rhs} of the previous equation.
We assume that this term renormalizes strongly the electronic band gap, but has no effect in the curvature of the bands.
So we start from that term with $\bk = 0$.
We denote the exchange energy by $\Delta_{ex}$ and it writes

\begin{equation}
    \Delta_{\mathrm{ex}} = \frac{1}{A} \sum_{\bp} V_{\bp} \left[F_{vvvv}(\bp ,0,\bp) - F_{vcvc}(\bp,0,\bp) \right] \,.
\end{equation}

The first form factor is

\begin{equation}
    F_{vvvv}(\bp ,0,\bp) = \langle v,0 | v, \bp \rangle \langle v,\bp | v, 0 \rangle = |\langle v,0 | v, \bp \rangle|^2.
\end{equation}

\noindent The overlap of the spinors (whose expressions one can consult in Eqs.~\eqref{eq:valence_spinor} and~\eqref{eq:conduction_spinor} reads

\begin{equation}
    \langle v,0 | v, \bp \rangle = \cos \frac{\xi_p}{2} \,, \qq{with} \xi_p = \arctan \frac{2 \hbar v_F p}{E_g}\,.
\end{equation}

\noindent This gives for the form factor

\begin{equation}
    F_{vvvv}(\bp ,0,\bp) = \cos^2\frac{\xi_p}{2} \,.
\end{equation}

The second form factor is

\begin{equation}
    F_{vcvc}(\bp ,0,\bp) = \langle c,0 | v, \bp \rangle \langle v,\bp | c, 0 \rangle = |\langle c,0 | v, \bp \rangle|^2.
\end{equation}

\noindent The spinors overlap (up to a phase) is now

\begin{equation}
    \langle c,0 | v, \bp \rangle = \sin \frac{\xi_p}{2} \,,
\end{equation}

\noindent hence

\begin{equation}
    F_{vcvc}(\bp ,0,\bp) = \sin^2\frac{\xi_p}{2} \,.
\end{equation}

\noindent Now we take the difference between the form factors

\begin{equation}
    \begin{split}
        F_{vvvv}(\bp ,0,\bp) - F_{vcvc}(\bp,0,\bp) &= \cos^2\frac{\xi_p}{2} - \sin^2\frac{\xi_p}{2} = \cos \xi_p = \cos \left(  \arctan \frac{2 \hbar v_F p(\beta)}{J_0(\beta) E_g} \right) = \\
        &= \frac{1}{\sqrt{1 + \frac{p^2 \gamma^2(\theta)}{k_0^2}}} \,,
    \end{split}
\end{equation}

\noindent where we defined the characteristic wavenumber

\begin{equation}
    k_0 \equiv |J_0 (\beta)| \frac{E_g}{2 \hbar v_F} \,.
\end{equation}

Using the auxiliar function $\gamma(\theta) = \sqrt{\cos^2 \theta + J_0^2(\beta) \sin^2 \theta}$, the exchange energy now reads

\begin{equation}
    \Delta_{\mathrm{ex}} = \frac{1}{A} \sum_{\bp} V_{\bp}  \frac{1}{\sqrt{1 + \left(\frac{p \gamma(\theta)}{k_0}\right)^2}} = \frac{e}{8 \pi^2 \epsz \kappa} \int_0^{2\pi} \dd \theta \int_0^{\infty} \dd p  \frac{1}{ (1 + r_0 p/\kappa) \sqrt{1 + \left(\frac{p \gamma(\theta)}{k_0}\right)^2}} \,.
\end{equation}

\noindent where we converted the sum into an integral as $\sum_{\bp} \to A/(2 \pi)^2  \int \dd \bp$, and used for the Rytova--Keldysh potential

\begin{equation}
    V_{\bp} = \frac{e}{2 \epsz p (\kappa + r_0 p)} \,.
\end{equation}

\noindent being $\kappa$ the arithmetic average between the dielectric constants of the substrate and the cladding material.


\noindent To perform the final integral, we make the change of variable $p \gamma(\theta) / k_0 = \tan u$.
The differential becomes $ \dd p = k_0 \dd(\tan u)/\gamma(\theta) = k_0 \frac{1}{\cos^2 u}/\gamma(\theta) \dd u$, and the square root in the denominator transforms into $\sqrt{1 + p^2 \gamma^2(\theta)/k_0^2} = \sqrt{1 + \tan^2 u} = 1/\cos u$.
When $p=0$ we have $u=0$, and when $p = \infty$ we have $u=\pi/2$.
Furthermore, we introduce a new dimensionless parameter $\gamma \equiv r_0 k_0$, so that the integral becomes now

\begin{equation}
    k_0 \int_0^{2\pi} \frac{1}{\gamma(\theta) \sqrt{1 + \alpha^2(\theta)}} \ln \left( \frac{1 + \alpha(\theta) + \sqrt{1 + \alpha^2(\theta)}}{1 + \alpha(\theta) - \sqrt{1 + \alpha^2(\theta)}} \right) \dd \theta
\end{equation}

\noindent Finally, we arrive at the expression for the exchange energy

\begin{equation}
    \Delta_{\mathrm{ex}} = \frac{e k_0}{8 \pi^2 \epsz \kappa} \int_0^{2\pi} \frac{1}{\gamma(\theta) \sqrt{1 + \alpha^2(\theta)}} \ln \left( \frac{1 + \alpha(\theta) + \sqrt{1 + \alpha^2(\theta)}}{1 + \alpha(\theta) - \sqrt{1 + \alpha^2(\theta)}} \right) \dd \theta \,\,
\end{equation}

\noindent which can be evaluated numerically without any computational effort.

